\section{Introduction}\label{sec:introduction}

For a mixing subshift of finite type \((X_A^+,\sigma)\) with primitive
adjacency matrix \(A\) and Perron root \(\lambda=\rho(A)\), the
dynamical \(\zeta\)-function
\[
  \zeta_A(z)
  =\exp\!\left(\sum_{n\ge 1}\frac{\Tr(A^n)}{n}z^n\right)
  =\det(I-zA)^{-1}
\]
has a simple pole at \(z=\lambda^{-1}\). The purpose of this paper is
to study the finite part at that pole and to show that it organises
both cyclic reconstruction and finite-group extensions.

After removing the Perron factor one is led to two closely related
quantities:
\[
  \FP(A)
  :=\lim_{r\uparrow 1}\bigl(P_A(r)+\log(1-r)\bigr),
  \qquad
  C(A):=\det\nolimits'(I-A/\lambda)^{-1}.
\]
Our first theme is that \(\FP(A)\) is not a residual correction term:
it is a primitive-orbit invariant with a closed formula in terms of
\(C(A)\) and the primitive orbit counts \(p_n(A)\).

The second theme is reconstruction. Writing
\[
  F_A(t):=\frac{\det(I-tA/\lambda)}{1-t},
\]
we show that the cyclic-lift constants
\(\{C(A^{\langle q\rangle})\}_{q\ge 1}\), and even one sufficiently
large root-of-unity layer, determine \(F_A\) and hence the full
non-Perron spectrum with multiplicity.

The third theme is non-abelian. For a finite-group extension defined by
an edge cocycle \(\tau:E_A\to G\), we introduce class-function
logarithms and show that primitive extraction is governed not by
M\"obius inversion alone, but by M\"obius inversion intertwined with
the Adams operations \(g\mapsto g^m\). This yields a primitive
determinant calculus in Peter--Weyl coordinates. Quotient
functoriality then identifies cyclic data as the exact abelian shadow
of the full primitive class profile, with an orthogonal genuinely
non-abelian defect.

\subsection{Main results}

The first flagship theorem identifies the finite part at the Perron
pole.

\begin{theorem}[Finite-part formula at the Perron pole]
  \label{thm:intro-finite-part}
  Let \(A\) be a primitive non-negative matrix with Perron root
  \(\lambda=\rho(A)>1\), reduced determinant \(C(A)\), and primitive
  orbit counts \(p_n(A)\). Then
  \[
    \FP(A)
    =\log C(A)
      +\sum_{n\ge 1}p_n(A)
        \bigl(\lambda^{-n}+\log(1-\lambda^{-n})\bigr).
  \]
  Equivalently,
  \[
    \FP(A)
    =\log C(A)+\sum_{k\ge 2}\frac{\mu(k)}{k}
      \log\det(I-\lambda^{-k}A)^{-1}.
  \]
  In particular, the orbit-Mertens constant
  \(\mathcal M(A)=\gamma+\FP(A)\) is determined by the reduced
  determinant and the primitive orbit counts.
\end{theorem}

This is proved in
\Cref{thm:finite-part-primitive-closed-form,cor:finite-part-mertens-asymptotic}.

The second flagship theorem shows that the cyclic side recovers the
reduced spectrum.

\begin{theorem}[Cyclic reconstruction of the reduced spectrum]
  \label{thm:intro-cyclic-reconstruction}
  Let \(A\) be a primitive non-negative \(d\times d\) matrix, let
  \[
    F_A(t)=\frac{\det(I-tA/\lambda)}{1-t},
    \qquad
    \Psi_q(A):=\log q+\log C(A^{\langle q\rangle}),
  \]
  and let \(\mu_2,\dots,\mu_d\) be the non-Perron eigenvalues of \(A\).
  Then:
  \begin{enumerate}[label=\textup{(\roman*)}]
    \item the sequence \(\{\Psi_q(A)\}_{q\ge 1}\) determines \(F_A\);
      equivalently, so does the family of cyclic-lift constants
      \(\{C(A^{\langle q\rangle})\}_{q\ge 1}\). Hence the cyclic data
      recover the non-Perron spectrum of \(A\) with algebraic
      multiplicity;
    \item if \(q\ge d\), then the single root-of-unity layer
      \(\{F_A(\omega_q^j)\}_{0\le j\le q-1}\) already determines every
      coefficient of \(F_A\);
    \item the exponential fingerprint satisfies
      \[
        \limsup_{q\to\infty}\abs{\Psi_q(A)}^{1/q}
        =
        \Lambda(A)/\lambda,
      \]
      where \(\Lambda(A)\) is the maximal modulus of the non-Perron
      spectrum.
  \end{enumerate}
\end{theorem}

This is the content of
\Cref{thm:finite-part-cyclic-lift-spectrum-identifiability,thm:finite-part-cyclic-lift-single-q-torsion-reconstruction,cor:finite-part-cyclic-lift-zeta-torsion-determines-spectrum}.

The finite-group side is the genuinely new non-abelian package.

\begin{theorem}[Primitive determinant calculus for finite-group
  extensions]
  \label{thm:intro-finite-group}
  Let \(\tau:E_A\to G\) be a cocycle into a finite group, let
  \(A_\varrho\) be the twisted transition matrices, and for a class
  function \(\phi:G\to\CC\) write
  \[
    \mathcal L_\phi(z)
    :=\sum_{n\ge 1}\frac{T_{\phi,n}}{n}z^n,
    \qquad
    \Pi_\phi(z):=\sum_{\gamma\in\cP}\phi(\tau(\gamma))z^{|\gamma|}.
  \]
  Then:
  \begin{enumerate}[label=\textup{(\roman*)}]
    \item class-function logarithms linearise the twisted determinant
      data, and Adams--M\"obius inversion expresses each primitive
      Peter--Weyl series directly in terms of the twisted determinants;
    \item if
      \(
        \eta:=\max_{\varrho\neq\mathbf 1}\rho(A_\varrho)<\lambda,
      \)
      then the conjugacy-class constants
      \[
        \Delta_C(A)
        :=
        \mathcal M_C(A)-\frac{|C|}{|G|}\mathcal M(A)
      \]
      reconstruct the non-trivial values
      \(\Pi_\varrho(\lambda^{-1})\) by non-abelian Fourier inversion;
    \item quotient maps \(G\twoheadrightarrow Q\) intertwine primitive
      class profiles. In particular, cyclic quotient observables span
      exactly the one-dimensional Peter--Weyl sector, so they exhaust
      the abelian shadow and are orthogonal to the genuine non-abelian
      defect.
  \end{enumerate}
  An explicit \(S_3\)-extension shows that vanishing periodic traces
  does not force vanishing primitive bias.
\end{theorem}

The detailed statements are proved in
\Cref{thm:class-function-linearisation,cor:class-function-adams-mobius,cor:primitive-peter-weyl-determinant,thm:class-mertens-explicit,thm:class-mertens-fourier,thm:quotient-functoriality,thm:abelian-shadow-defect,thm:cyclic-detection-boundary,cor:s3-example-profile}.

\subsection{What is standard and what is new}

The background used here is classical: rationality of
\(\zeta\)-functions for shifts of finite type
\cite{BowenLanford1970Zeta,Manning1971Axiom}, the transfer-operator
viewpoint of Ruelle and Parry--Pollicott
\cite{Ruelle1976ZetaExpanding,Ruelle1978Thermodynamic,ParryPollicott1990Zeta},
and the twisted determinant formalism for finite-group extensions and
periodic data
\cite{ParryPollicott1986Chebotarev,NooraniParry1992ChebotarevShifts,PollicottSharp2007Chebotarev,BoyleSchmieding2017FiniteGroupExtensions,Parry1999LivsicNonAbelian}.
For graph coverings one should also compare Ihara--Bass type
determinantal identities and their weighted variants
\cite{Bass1992Ihara,ChoeKwakParkSato2007Bartholdi,Sato2008WeightedBartholdiCoverings}.

What is not standard is the way the primitive data are organised. On
the abelian side, the finite part at the Perron pole becomes a
reconstruction invariant. On the non-abelian side, the primitive
series are not recovered by a bare M\"obius inversion: Adams
operations are forced into the determinant calculus, and quotient
functoriality identifies the precise frontier between cyclic detection
and genuinely non-abelian information.

\begin{center}
\small
\begin{tabular}{p{0.18\textwidth}p{0.20\textwidth}p{0.24\textwidth}p{0.28\textwidth}}
\toprule
Present theorem & Closest antecedent & Standard input used here & New theorem-level input \\
\midrule
Finite-part formula (\Cref{thm:finite-part-primitive-closed-form}) &
Euler products and zeta rationality
\cite{BowenLanford1970Zeta,Manning1971Axiom,ParryPollicott1990Zeta} &
Primitive-orbit expansion of \(\zeta_A\), Perron factorisation of
\(\det(I-zA)\) &
The finite part \(\FP(A)\) is written as an explicit primitive-orbit
invariant tied to the reduced determinant \(C(A)\). \\

Cyclic reconstruction
(\Cref{thm:finite-part-cyclic-lift-spectrum-identifiability,thm:finite-part-cyclic-lift-single-q-torsion-reconstruction}) &
Determinant formalism for cyclic covers and root-of-unity sampling &
Spectra of cyclic lifts, discrete Fourier inversion on roots of unity &
Cyclic-lift constants reconstruct the reduced spectral polynomial, and
one sufficiently large root-of-unity layer already suffices. \\

Primitive determinant calculus
(\Cref{thm:class-function-linearisation,cor:class-function-adams-mobius,cor:primitive-peter-weyl-determinant}) &
Twisted determinant factorisations and Peter--Weyl expansions for
finite-group extensions &
Character expansions, determinant factorisation, Adams operations on
class functions &
Primitive extraction is expressed directly in determinant form via an
Adams--M\"obius inversion, producing explicit primitive
Peter--Weyl series. \\

Class constants, quotient shadows, and cyclic detection boundary
(\Cref{thm:class-mertens-fourier,thm:quotient-functoriality,thm:cyclic-detection-boundary}) &
Chebotarev-type periodic counting and quotient constructions for
finite-group extensions &
Fourier inversion on class functions, functoriality under quotient
maps &
The conjugacy-class constants become exact non-abelian Fourier
coordinates of primitive data, and cyclic quotients are shown to
exhaust precisely the abelian shadow. \\
\bottomrule
\end{tabular}
\end{center}

Accordingly, the finite-group part of the paper is not another
Chebotarev statement in disguise. Its novelty lies in a primitive-level
determinant calculus and in the exact identification of the information
that survives passage to cyclic quotients.

\subsection{Organisation}

\Cref{sec:preliminaries} fixes notation for primitive orbit counts,
reduced determinants, cyclic lifts, and finite-group extensions, and
separates the purely matrix-theoretic part of the paper from the later
graph-extension setting. \Cref{sec:finite-part} treats the Perron-pole
finite part and the cyclic reconstruction of the reduced spectrum.
\Cref{sec:chebotarev} develops the primitive determinant calculus for
finite-group extensions, isolates the genuinely new Adams--M\"obius
mechanism, and derives the class constants under a twisted spectral
gap. \Cref{sec:shadows} proves quotient functoriality, identifies the
abelian shadow and the exact boundary of cyclic detection, and closes
with the explicit \(S_3\)-model. The appendices collect holomorphic
variation, further root-of-unity consequences, and the rigidity results
for growing families.
